\PassOptionsToPackage{dvipsnames,svgnames,table}{xcolor}
\documentclass[11pt]{article}
\usepackage[a4paper,margin=22mm]{geometry}
\usepackage[no-math]{fontspec}
\setmainfont{DejaVu Serif}
\usepackage{amsmath,amssymb,amsthm,mathtools,graphicx,xcolor,hyperref,xurl,xspace,ifthen,enumerate}
\usepackage[shortlabels]{enumitem}
\setlength{\emergencystretch}{4em}
\AtBeginDocument{\special{pdf:mapline pzdr Dingbats <uzdr.pfb}}
\SetMathAlphabet{\mathbf}{normal}{OT1}{cmr}{bx}{n}
\SetMathAlphabet{\mathbf}{bold}{OT1}{cmr}{bx}{n}
\newtheorem{theo}{Theorem}[section]
\newtheorem{lemm}[theo]{Lemma}
\newtheorem{thm}[theo]{Theorem}
\newtheorem{proposition}[theo]{Proposition}
\newtheorem{corollary}[theo]{Corollary}
\newtheorem{definition}[theo]{Definition}
\newtheorem{theorem}[theo]{Theorem}
\newtheorem*{remas*}{Remarks}
\newtheorem{remas}[theo]{Remarks}
\newtheorem*{exam*}{Example}
\newtheorem{ques}[theo]{Question}
\providecommand{\og}{«\nobreakspace}
\providecommand{\fg}{\nobreakspace»}
\newtheorem{lem}[theo]{Lemma}
\newtheorem{lemma}[theo]{Lemma}
\newtheorem{prop}[theo]{Proposition}
\newtheorem{coro}[theo]{Corollary}
\newtheorem{corol}[theo]{Corollary}
\newtheorem*{conj*}{Conjecture}
\newtheorem{conj}[theo]{Conjecture}
\newtheorem{conjecture}[theo]{Conjecture}
\newtheorem{Proposition}[theo]{Proposition}
\newtheorem{theoreme}[theo]{Theorem}
\theoremstyle{definition}
\newtheorem{defi}[theo]{Definition}
\newtheorem{exam}[theo]{Example}
\newtheorem{exem}[theo]{Example}
\newtheorem{remark}[theo]{Remark}
\newtheorem{example}[theo]{Example}
\newtheorem{exams}[theo]{Examples}
\newtheorem{rema}[theo]{Remark}
\newtheorem{nota}[theo]{Notation}
\newtheorem{notation}[theo]{Notation}
\newtheorem*{notas}{Notations}
\newtheorem{result}[theo]{Result}
\newtheorem{question}[theo]{Question}
\newtheorem{prob}[theo]{Problem}
\newtheorem*{rema*}{Remark}
\newtheorem*{defi*}{Definition}
\newtheorem*{theo*}{Theorem}
\newtheorem*{lemm*}{Lemma}
\newtheorem*{prop*}{Proposition}
\newcommand{\enoncename}{}
\newtheorem{corpusenonce}[theo]{\enoncename}
\newenvironment{enonce}[2][]{\renewcommand{\enoncename}{#2}\begin{corpusenonce}}{\end{corpusenonce}}
\newenvironment{enonce*}[2][]{\par\medskip\noindent\textbf{#2.}\itshape}{\par\medskip}
\newenvironment{proofc}[1][\proofname]{\begin{proof}[#1]}{\end{proof}}
\newcommand{\Mc}{\mathcal{M}}
\newcommand{\pointir}{.\ ---\ }
\newcommand{\equalenv}[2]{\expandafter\let\csname#1\expandafter\endcsname\csname#2\endcsname\expandafter\let\csname end#1\expandafter\endcsname\csname end#2\endcsname}
\providecommand{\diagbox}[3][]{\begin{tikzpicture}[x=1em,y=1em,baseline=(current bounding box.center)]\useasboundingbox(0,0) rectangle(3,1.4);\draw(0,1.4)--(3,0);\node at(.5,.3){#2};\node at(2.5,1.1){#3};\end{tikzpicture}}
\providecommand{\eg}{e.g.\ }
\providecommand{\ie}{i.e.\ }
\providecommand{\cf}{cf.\ }
\providecommand{\longthanks}[1]{\par\smallskip\noindent #1\par\smallskip}
\providecommand{\qbinom}[2]{\genfrac{[}{]}{0pt}{}{#1}{#2}}

\providecommand{\mathof}[1]{\mathrm{#1}}

\numberwithin{equation}{section}


\def\endappendix{}


\usepackage{latexsym,amsmath,amssymb,verbatim,tikz,tikz-cd}
\usepackage{mathtools}
\usepackage{graphicx,xcolor}

\usepackage[english]{babel}
\usepackage{amssymb}
\usepackage{amsmath}
\usepackage{amsthm}
\usepackage{graphicx}
\usepackage{csquotes}
\usepackage{bm}
\usepackage{enumitem}
\usepackage{subcaption}
\usepackage{caption}
\renewcommand{\thesubfigure}{\Alph{subfigure}}
\usepackage{dirtytalk}
\providecommand{\nopunct}{\spacefactor1007\relax}
\newtheorem{problem}[theo]{Problem}

\newcommand{\sym}{\mathrm{Sym}}
\newcommand{\qsym}{\mathrm{QSym}}
\newcommand{\Des}{\operatorname{Des}}
\newcommand{\Asc}{\operatorname{Asc}}
\newcommand{\Syt}{\operatorname{SYT}}
\newcommand{\Short}{\operatorname{Short}}
\newcommand{\Loop}{\operatorname{Loop}}
\newcommand{\Stable}{\operatorname{Stable}}
\newcommand{\len}{\operatorname{len}}
\newcommand{\height}{\operatorname{height}}
\newcommand{\core}{\operatorname{core}}
\newcommand{\res}{\operatorname{res}}
\newcommand{\row}{\operatorname{row}}
\newcommand{\insrt}{\operatorname{insert}}
\newcommand{\UP}{\operatorname{UP}}
\newcommand{\DOWN}{\operatorname{DOWN}}


\newcommand{\ZZ}{\mathbb{Z}}
\newcommand{\NN}{\mathbb{N}}
\newcommand{\RR}{\mathbb{R}}
\newcommand{\QQ}{\mathbb{Q}}
\newcommand{\PP}{\mathbb{P}}

\newcommand{\F}{\mathcal{F}}
\newcommand{\Q}{\mathcal{Q}}
\newcommand{\A}{\mathcal{A}}
\newcommand{\M}{\mathcal{M}}
\newcommand{\I}{\mathcal{I}}
\newcommand{\PM}{\mathcal{PM}}
\newcommand{\Pcal}{\mathcal{P}}


\newcommand{\odds}{\mathrm{odds}}

\makeatletter
\newcommand{\proofstep}[1]{
  \par
  \addvspace{\medskipamount}
  \textbf{#1\@addpunct{.}}\enspace\ignorespaces
}

\newcommand\restr[2]{{
  \left.\kern-\nulldelimiterspace 
  #1 
  \vphantom{\big|} 
  \right|_{#2} 
  }}
\makeatother












\title{The descent set of the tableau of a matching is its set of short chords}
\author{Avichai Marmor}
\date{doi:10.5802/alco.351; CC BY 4.0}
\begin{document}
\maketitle
\noindent\textit{Editorial note.} The counted statement is the author's labelled proposition environment at publisher-source.tex lines 897-900; its proof environment opens at line 901, immediately after \textbackslash{}end\{proposition\}, so the pairing is by adjacency and the retained proof range 901-906 is the whole environment.\par\medskip
\noindent\textit{Editorial note.} The lemma lem:properties of core, with its three numbered items (source lines 873-881), is retained as context because the counted proof cites its first item at source line 905; it is the only outside label of the retained proof.\par\medskip
\noindent\textit{Editorial note.} This article carries no label-number dump and the retained counter values are left empty for the whole source, so the printed numbering of the card is the wrapper's default; the identity of the counted item is its author label prop:F statistic preserving.\par\medskip
\noindent\textit{Editorial note.} r002 note: r001 failed with an undefined reference to def:restriction of matching, which the retained lem:properties of core context cites at source line 879; r002 adds exactly that definition environment as a further context and changes nothing else.\par\medskip
\noindent\textit{Editorial scope.} One original result of the pinned publisher article, counted once: the descent set of the tableau of a matching is its set of short chords (author label prop:F statistic preserving).\par\medskip
\noindent\textit{Allowed standard prerequisites (editorial).}\begin{itemize}
\item Core and stable vertices. The core of a matching is the matching induced on its stable vertices and the two-row tableau T(m) records the unstable chords in its second row.. They define the tableau whose descent set is computed.
\item Short chord and descent set. A chord (i,i+1) is short when no vertex lies strictly between i and i+1, and the descent set of a tableau is the set of indices i with i+1 in the row below i.. They are the two sets that the proposition identifies.
\item Stability of a chord intersecting a non-crossing chord. A chord that intersects another chord is stable.. It is the contradiction used in the proof of the reverse inclusion.
\end{itemize}\medskip
\begin{definition}
\label{def:restriction of matching}
    Given a matching $m \in \M_N$, we say that a set $S \subseteq [N]$ is $m$-invariant if for every chord $(i,j) \in m$ we have $i \in S$ if and only if $j \in S$. Moreover, for a matching $m$ and an $m$-invariant set $S$ we define the \emph{restriction} of $m$ to $S$, denoted $\res_S(m)$, to be a matching on $S$ which is obtained by removing all the vertices not in $S$ from $m$.
\end{definition}

\begin{lemma}
\label{lem:properties of core}
    The reduction process has the following properties:
    \begin{enumerate}[label=(\arabic*), ref=\ref{lem:properties of core}\text{ (part~}\arabic*)]
        \item \label{lem:properties of core--intersect stable} If a chord intersects another chord then it is stable.
        \item \label{lem:properties of core--chord unstable iff all midpoints unstable} A chord $(i,j)$ is stable if and only if there exists a stable vertex in $\{i+1,\dots,j-1\}$.
        \item \label{lem:properties of core--perfect non-crossing unstable} Given a matching $m$ and $i < j$, if the set $[i,j]$ is $m$-invariant \textup{(}as defined in Definition~\ref{def:restriction of matching}\textup{)} and the restricted matching $\res_{[i,j]}(m)$ is perfect and non-crossing, then $\ell \notin \Stable(m)$ for all $\ell \in [i,j]$.
    \end{enumerate}
\end{lemma}

\begin{proposition}
\label{prop:F statistic preserving}
    Let $m \in \M_{N,f}$ be a matching, and denote $T := T(m)$. Then $\Des T = \Short m$.
\end{proposition}

\begin{proof}
    Let $i \in \Short m$ be an index representing a short chord $(i,i+1) \in m$. The chord $(i,i+1)$ is unstable, so we may deduce from the definition of $T(m)$ that $i \in \row_1(T)$ and $i+1 \in \row_2(T)$. Therefore, $i \in \Des T$.
    
    On the other hand, let $i \in \Des T$ be a descent of $T$. This implies that $i \in \row_1(T)$ and $i+1 \in \row_2(T)$. Since $i+1 \in \row_2(T)$, we obtain that there exists $j < i+1$ such that $(j,i+1) \in m$ is an unstable chord. Since the chord $(j,i+1)$ is unstable and $j \le i \le i+1$, we obtain by Lemma~\ref{lem:properties of core--chord unstable iff all midpoints unstable} that $i$ is an endpoint of an unstable chord of $m$. Since $i \in \row_1(T)$, we obtain that $i$ opens an unstable chord $(i,j')$ of $m$ for some $j'>i$.
    If $(i,i+1) \notin m$ then $j<i$ and $j' > i+1$, and we obtain that the unstable chord $(j,i+1)$ intersects the chord $(i,j')$, contradicting Lemma~\ref{lem:properties of core--intersect stable}. Therefore, we may conclude that $(i,i+1) \in m$ and $i \in \Short m$.
\end{proof}

\end{document}
